\documentclass{beamer}

% ============================================================
% THEME & COLORS
% ============================================================
\usetheme{Madrid}
\usecolortheme{whale}

% Custom color scheme (Premium Blue style)
\definecolor{customBlue}{RGB}{29, 78, 216}     % Royal Blue
\definecolor{customDarkBlue}{RGB}{30, 58, 138} % Deep Navy Blue
\setbeamercolor{palette primary}{bg=customBlue, fg=white}
\setbeamercolor{palette secondary}{bg=customDarkBlue, fg=white}
\setbeamercolor{palette tertiary}{bg=black, fg=white}
\setbeamercolor{titlelike}{parent=palette primary}
\setbeamercolor{structure}{fg=customBlue}

% Packages
\usepackage[utf8]{inputenc}
\usepackage[vietnamese]{babel}
\usepackage{booktabs}

% ============================================================
% TITLE PAGE
% ============================================================
\title[myShell]{Báo cáo Bài tập lớn Hệ điều hành\\Chương trình Shell đơn giản (myShell)}
\author[Phạm Văn Sâm]{Sinh viên thực hiện: Phạm Văn Sâm\\MSSV: 202400114}
\institute[HUST]{Trường Đại học Bách Khoa Hà Nội}
\date{\today}

\begin{document}

% Slide 1: Title
\begin{frame}
    \titlepage
\end{frame}

% ============================================================
% PHẦN 1: GIỚI THIỆU KHÁI NIỆM & DANH SÁCH LỆNH
% ============================================================
\section{Khái niệm \& Danh sách lệnh}

\begin{frame}{Shell là gì? Chu trình REPL}
    \begin{block}{Khái niệm đơn giản}
        **Shell** là phần mềm giao diện dòng lệnh giúp người dùng tương tác trực tiếp với nhân hệ điều hành (Kernel).
    \end{block}
    \begin{block}{Chu trình hoạt động (REPL)}
        \begin{itemize}
            \item **Read:** Đọc dòng lệnh người dùng gõ từ bàn phím.
            \item **Eval:** Phân tích cú pháp và thực thi lệnh (Built-in hoặc tạo tiến trình con qua `CreateProcess`).
            \item **Print:** In kết quả ra màn hình.
            \item **Loop:** Lặp lại chu trình vô hạn.
        \end{itemize}
    \end{block}
\end{frame}

\begin{frame}{Danh sách lệnh nội trú (Built-in Commands)}
    Các lệnh được thực thi trực tiếp ngay trong tiến trình của Shell:
    \begin{itemize}
        \item \textbf{`help`}: Hiển thị bảng hướng dẫn sử dụng shell.
        \item \textbf{`exit`}: Giải phóng các tài nguyên và thoát khỏi shell.
        \item \textbf{`cd <path>`}: Chuyển thư mục làm việc hiện hành. Gõ `cd` không tham số để in đường dẫn hiện hành.
        \item \textbf{`dir [path]`}: Liệt kê nội dung thư mục hiện tại hoặc thư mục chỉ định.
        \item \textbf{`date`}: Xem ngày hiện tại của hệ thống.
        \item \textbf{`time`}: Xem giờ hiện tại của hệ thống.
    \end{itemize}
\end{frame}

\begin{frame}{Quản lý đường dẫn (Path Management)}
    Các lệnh thao tác với các đường dẫn tìm kiếm chương trình tùy chỉnh:
    \begin{itemize}
        \item \textbf{`path`}: Xem danh sách các thư mục tìm kiếm tùy chỉnh đang được cấu hình trong shell.
        \item \textbf{`addpath <path>`}: Thêm một thư mục vào danh sách tìm kiếm (Tự động lưu dạng đường dẫn tuyệt đối và đồng bộ vào biến môi trường `PATH` của tiến trình).
        \item \textbf{`delpath <path>`}: Xóa đường dẫn khỏi danh sách tìm kiếm bằng tên trực tiếp (không cần ID/index) và cập nhật lại biến môi trường `PATH`.
    \end{itemize}
\end{frame}

\begin{frame}{Quản lý tiến trình (Process Management)}
    Chạy chương trình ngoài ở hai chế độ:
    \begin{itemize}
        \item **Foreground:** Chạy trực tiếp, shell sẽ đợi tiến trình con chạy xong.
        \item **Background:** Thêm dấu `&` ở cuối lệnh. Shell chạy chương trình ngầm và trả lại dấu nhắc lệnh ngay lập tức.
    \end{itemize}
    \begin{block}{Các lệnh điều khiển tiến trình chạy ẩn}
        \begin{itemize}
            \item \textbf{`list`}: Xem danh sách các tiến trình chạy nền kèm PID và trạng thái.
            \item \textbf{`stop <pid>`}: Tạm dừng một tiến trình nền đang chạy.
            \item \textbf{`resume <pid>`}: Tiếp tục chạy tiến trình nền đang bị tạm dừng.
            \item \textbf{`kill <pid>`}: Dừng cưỡng bức (hủy) một tiến trình nền.
        \end{itemize}
    \end{block}
\end{frame}

\begin{frame}{Xử lý Ctrl+C \& Tệp batch (`.bat`)}
    \begin{block}{Xử lý tín hiệu ngắt (Ctrl+C)}
        \begin{itemize}
            \item Khi bấm `Ctrl+C`, nếu có tiến trình Foreground đang chạy, tiến trình đó sẽ bị dừng lập tức.
            \item Chương trình Shell cha vẫn được giữ sống để người dùng tiếp tục thao tác.
        \end{itemize}
    \end{block}
    \begin{block}{Thực thi tệp kịch bản `.bat`}
        \begin{itemize}
            \item Đọc và chạy lần lượt các dòng lệnh cấu hình trong tệp `.bat`.
            \item Hỗ trợ ẩn hiển thị lệnh (`@echo off`) và chạy dưới nền (ví dụ: `test.bat &`).
        \end{itemize}
    \end{block}
\end{frame}

% ============================================================
% PHẦN 2: KỊCH BẢN DEMO
% ============================================================
\section{Kịch bản Demo}

\begin{frame}{Kịch bản Demo 1: Lệnh nội trú \& Biến môi trường}
    Các bước gõ lệnh Demo trực tiếp cho giảng viên xem:
    \begin{enumerate}
        \item **In giờ ngày và trợ giúp:**
              \begin{itemize}
                  \item Gõ `help`
                  \item Gõ `date`, `time`
              \end{itemize}
        \item **Thao tác thư mục:**
              \begin{itemize}
                  \item Gõ `cd` (In thư mục hiện tại)
                  \item Gõ `dir` (Liệt kê danh sách file)
              \end{itemize}
        \item **Đồng bộ hóa PATH (addpath/delpath):**
              \begin{itemize}
                  \item Gõ `addpath test` (Thêm thư mục tuyệt đối)
                  \item Gõ `path` (Xem danh sách)
                  \item Di chuyển sang thư mục khác `cd src` rồi chạy thử tệp thực thi trong thư mục `test` bằng tên (ví dụ: `hello`).
                  \item Gõ `delpath C:\...\test` để xóa đường dẫn bằng tên.
              \end{itemize}
    \end{enumerate}
\end{frame}

\begin{frame}{Kịch bản Demo 2: Quản lý tiến trình \& Ctrl+C}
    \begin{enumerate}
        \item **Chạy Foreground và dừng bằng Ctrl+C:**
              \begin{itemize}
                  \item Gõ `ping -t 127.0.0.1` (Chương trình chạy liên tục)
                  \item Nhấn tổ hợp phím `Ctrl + C` $\rightarrow$ tiến trình dừng và Shell hiển thị lại dấu nhắc lệnh bình thường.
              \end{itemize}
        \item **Chạy Background và kiểm soát:**
              \begin{itemize}
                  \item Gõ `ping -t 127.0.0.1 &` $\rightarrow$ Tiến trình chạy ẩn dưới nền.
                  \item Gõ `list` $\rightarrow$ Xem PID và trạng thái tiến trình nền.
                  \item Gõ `stop <pid>` $\rightarrow$ Tạm ngưng tiến trình nền.
                  \item Gõ `list` $\rightarrow$ Trạng thái chuyển sang `Stopped`.
                  \item Gõ `resume <pid>` $\rightarrow$ Cho chạy tiếp.
                  \item Gõ `kill <pid>` $\rightarrow$ Hủy tiến trình hoàn toàn.
              \end{itemize}
    \end{enumerate}
\end{frame}

\begin{frame}{Kịch bản Demo 3: Thực thi tệp `.bat`}
    \begin{enumerate}
        \item **Chạy tệp script tuần tự:**
              \begin{itemize}
                  \item Gõ `test/2.bat` $\rightarrow$ Shell đọc và in kết quả chạy của từng dòng lệnh bên trong tệp batch.
              \end{itemize}
        \item **Chạy tệp script dưới nền (Background):**
              \begin{itemize}
                  \item Gõ `test/2.bat &` $\rightarrow$ Chạy ngầm tệp kịch bản.
                  \item Gõ `list` để xem trạng thái tiến trình nền thực thi script.
              \end{itemize}
    \end{enumerate}
\end{frame}

\begin{frame}
    \begin{center}
        \Huge \color{customBlue} Cảm ơn thầy đã theo dõi!\\
        \vspace{1cm}
        \Large Câu hỏi \& Đóng góp ý kiến
    \end{center}
\end{frame}

\end{document}
